\section{Introduction}
\label{sec:introduction}

The most common objective when training text embeddings is to place texts with similar semantics close together in the embedding space \citep{sentencetransformers}. Style representations, by contrast, aim to embed texts with similar writing styles near each other and texts with different styles far apart, regardless of their semantic content \citep{styleemb}. Embeddings are usually trained via contrastive learning, with triplets consisting of an anchor text, a positive text (which should be embedded closely to the anchor), and a negative text (which should be embedded far from the anchor) \citep{contrastive1, contrastive2, tripletloss}. Existing approaches often use social media datasets with the assumption that all writing by the same author shares a similar style and that texts by different authors exhibit dissimilar styles \citep{styleemb}. These methods also attempt to minimize content representation in the resulting embeddings. They select a text from the same author on a \textit{different topic} as the positive example, and a text from a different author on the \textit{same topic} as the negative example, approximating topic similarity using subreddit or conversation metadata. However, these methods are limited by the imperfect nature of data acquired under such assumptions. For example, the same author may write about the same topic even in different subreddits. As a result, these imperfect contrastive triplets do not explicitly control for content, leading to style embeddings with weak content-independence.
%\textcolor{gray}{A solution to this problem would be to use parallel datasets where style varies while content is held constant, but such datasets are rare. Without parallel data, training must instead rely on proxy objectives, such as using authorship-based contrastive triplets, which can introduce content leakage into style representations (see Figure \ref{figure:wegmannvsstyledistance}).} 
The ``content leakage'' caused by such proxy objectives (illustrated in Figure \ref{figure:wegmannvsstyledistance}) can undermine the effectiveness of style representations in tasks that require strict separation between style and content, such as stylistic analysis, authorship tasks, style transfer steering, and automatic style transfer evaluation. To overcome this, a more controlled approach to style contrastive learning is necessary.

In this paper, we introduce \textsc{StyleDistance}, a novel method for training stronger, content-independent style embeddings which leverages synthetic parallel text examples generated by a large language model (LLM) \citep{gpt4}. By creating near-exact paraphrases with controlled stylistic variations, we produce positive and negative examples across 40 distinct style features. This synthetic dataset, which we call \textsc{SynthStel}, enables more precise contrastive learning (visualized in Figure \ref{fig:main}) and is more robust to the content leakage inherent in existing datasets. We evaluate our method on both human and automated benchmarks, measuring the content-independence, quality, and utility of \textsc{StyleDistance} embeddings.

In summary, our primary contributions are:

\begin{enumerate}
\item We generate and release \textsc{SynthStel}, a dataset of near-exact paraphrases across 40 distinct style features.

\item We introduce \textsc{StyleDistance}, a new approach to style representation learning which uses synthetic parallel examples with controlled stylistic variations. We release the embedding model as a resource.

\item We demonstrate that \textsc{StyleDistance} significantly improves the content-independence of style embeddings, generalizing effectively to real-world benchmarks of style representation quality and outperforming existing style representations in downstream applications.
\end{enumerate}